\section{Experimental Setup}

We validate the frugal flow (FF) on MorphoMNIST under a family of synthetic interventions
designed so that the target estimand is known \emph{exactly} --- imposed by construction rather
than approximated by Monte Carlo. The experiments form a ladder in which one design feature
changes at a time: randomised assignment $\to$ confounded assignment $\to$ each of several
distinct kinds of effect variation. This isolates the contribution of each difficulty to the
estimation error.

All datasets are produced by a single parameterised generator
(\texttt{prepare\_morphomnist\_exps.py}) and all fits by a single estimation script
(\texttt{exp\_ate\_recovery.py}); the experiments below differ only in the value of a handful
of scalar knobs.

\subsection{Outcome representation}

The outcome is the image itself. MorphoMNIST images ($28 \times 28$ greyscale, 60{,}000
training images) are padded to $32 \times 32$, normalised to $[-1, 1]$, and downsampled by
adaptive average pooling to $s \times s$, giving an outcome of dimension $K = s^2$. We report
results at $s = 8$ ($K = 64$); $s$ is a free parameter, trading spatial resolution against the
dimensionality the FF must model jointly.

Pixel intensities are bounded and discretely quantised, so they are mapped to the unbounded
support a continuous flow requires. Writing $x \in [0,1]$ for the rescaled intensity, we apply
uniform dequantisation at the original 256-level granularity, squeeze into
$(\alpha/2,\, 1 - \alpha/2)$ with $\alpha = 0.05$ to avoid divergence at the endpoints, and
take the logit:
\begin{equation}
    Y \;=\; \operatorname{logit}\!\big(\alpha/2 + (1-\alpha)\,\tilde{x}\big),
    \qquad
    \tilde{x} = x + \varepsilon, \quad \varepsilon \sim \mathrm{U}[0, 1/256].
    \label{eq:logit-transform}
\end{equation}
The map is invertible, so samples can be pushed back to pixel space for visualisation.

\textbf{All estimands in this section live in logit space}, the space in which the FF is
fitted and in which the treatment effect is additive by construction. Because
$\operatorname{logit}^{-1}$ is nonlinear, it is not estimand-preserving: pixel-space
differences of means are \emph{not} a valid comparison target, and are used for figures only.

\subsection{The intervention}

Treatment applies an additive, spatially structured shift in logit space. For unit $i$ and
pixel $k$,
\begin{equation}
    Y_i(0) = Y_i, \qquad Y_i(1)_k = Y_i(0)_k + \tau_{ik},
    \qquad
    \tau_{ik} = m_k \cdot c_{ik},
    \label{eq:potential-outcomes}
\end{equation}
where $m \in \R^K$ is a fixed \emph{effect map} and $c_{ik}$ a per-unit \emph{modulation}
specified in Section~\ref{sec:effect-construction}. The baseline map is a centred disc of
radius $r = 2$ pixels (12 of the 64 pixels at $s = 8$) with magnitude $\delta = 1.0$ in logit
units; a ring, a constant map, and a smooth horizontal gradient are also available, the last
being the hardest case since it has neither flat regions nor exact zeros for an estimator to
exploit. The radius defaults to $\lceil s/4 \rfloor$, holding the map at $\approx 20\%$ of
pixels as $K$ varies.

\subsection{Pretreatment covariates}

The observed covariate vector $Z$ is mapped to uniform ranks $U_Z$ by the FF's stage-one
marginal transforms.

\begin{itemize}
    \item \textbf{Stroke thickness} $A \in \R$ (continuous), supplied with MorphoMNIST as the
    mean of the distance transform over the foreground skeleton. Thickness drives treatment
    assignment throughout, and modulates the effect in the heterogeneous experiments. Mapped
    to a rank $U_A$ by a univariate marginal NSF.

    \item \textbf{Image brightness} $B \in \R$ (continuous), the MorphoMNIST intensity
    summary. Used as a second effect modifier where stated, and included in $Z$ exactly when
    the effect depends on it --- a covariate the effect uses must be observed, or the
    conditional effect is not identified.

    \item \textbf{Digit class} $C \in \{0, \ldots, 9\}$ (discrete). Stroke morphology varies
    systematically across classes, so digit class is a natural additional covariate. One-hot
    encoded and mapped to a rank $U_C$ via the empirical CDF and the distributional transform.
\end{itemize}

The headline configuration restricts to a single digit class, so $Z$ contains only the
continuous block and the discrete stage is bypassed; this removes class variation as a
nuisance source and makes the confounding attributable to fully observed variables. An
all-classes variant ($n = 60{,}000$, $\dim Z = 11$) is generated by the same code path and
exercises the mixed continuous/discrete stage-one transform.

\subsection{Causal status of the covariates}
\label{sec:covariate-status}

Because the outcome here is an image and the covariates are morphometric quantities, it is
worth stating precisely why $A$ and $B$ are pretreatment covariates rather than functions of
the outcome.

The images are the synthetic MorphoMNIST dataset of \citet{pawlowski2020deepscm}, whose
structural causal model is known in closed form. Attributes are sampled first and the image is
\emph{synthesised to realise them}:
\begin{align}
    A &= 0.5 + \mathrm{Gamma}(10, 5), \label{eq:scm-thickness}\\
    B &= 191\,\sigma\!\big(\mathcal{N}((A - 2.5)\cdot 2,\; 0.5)\big) + 64,
        \label{eq:scm-intensity}\\
    x &= \operatorname{clip}\!\Big(\mathrm{SetThickness}(A)\big(\xi\big)
         \cdot B \,/\, \overline{B}(\cdot),\; 0,\; 255\Big),
        \label{eq:scm-image}
\end{align}
where $\xi$ is the source MNIST digit, which plays the role of the image mechanism's exogenous
noise. We verified the released files against
Equations~\eqref{eq:scm-thickness}--\eqref{eq:scm-intensity}: the realised thickness has mean
$2.503$ and standard deviation $0.632$ against theoretical $2.500$ and $0.632$, and the
intensity residual on the logit scale has standard deviation $0.5007$ against the generator's
$0.5$. The attribute columns are therefore the \emph{sampled} values, not quantities
re-measured from the image, so there is no measurement error between the latent factors and
what we condition on.

Consequently $A$ and $B$ are causal \emph{parents} of the image, not summaries of it, and the
full graph is
\begin{equation}
    A \to B, \quad A \to x, \quad B \to x, \quad \xi \to x, \quad
    A \to T, \quad x = Y(0) \to Y(1),
\end{equation}
with $T$ generated by Equation~\eqref{eq:propensity} from $A$ and an independent uniform draw.
Three properties follow.

\begin{enumerate}
    \item \textbf{$Z$ is a non-descendant of $T$.} The attributes are generated before, and
    independently of, the assignment mechanism, so $Z^{do(T=1)} = Z^{do(T=0)} = Z$ --- the
    definition of a pretreatment covariate. This is verifiable rather than assumed: holding the
    seed fixed and varying $\beta_1$ from $0$ to $3$, which moves
    $\operatorname{corr}(A, T)$ from $-0.01$ to $+0.69$, leaves $A$, $B$ and $Y(0)$
    bitwise identical.

    \item \textbf{$Z$ is a confounder in the textbook sense}, a common cause of treatment and
    outcome: $A \to T$ through the propensity and $A \to x = Y(0)$ through the morphological
    operation. It is neither a mediator (as $T$ does not cause $A$ or $B$) nor a collider
    (as neither has an incoming edge from $T$ or from $Y$). The only back-door paths,
    $T \leftarrow A \to x \to Y$ and $T \leftarrow A \to B \to x \to Y$, are blocked by
    conditioning on $A$, so $\{A\}$ satisfies the back-door criterion.

    \item \textbf{Conditional ignorability holds by construction, not by assumption.} Since
    $T = \mathbb{1}\{U < \sigma(\beta_0 + \beta_1 \tilde{A})\}$ with $U$ drawn independently of
    every other variable, $T \perp (Y(0), Y(1)) \mid A$ holds exactly. Positivity likewise:
    propensities are bounded away from $0$ and $1$ and both extremes are reported per run. The
    oracle weighting result of Section~\ref{sec:floor} is the empirical confirmation --- an
    invalid adjustment set could not recover the imposed effect to within sampling noise.
\end{enumerate}

We stress that the recoverability of $A$ from $x$ --- a linear predictor on the $64$ pixels
attains $R^2 = 0.94$ --- is not evidence against this ordering. $\mathrm{SetThickness}$ writes
$A$ into the image, so the image necessarily reveals it; inverting a mechanism does not reverse
its arrow. The discriminating test is interventional, and it is the invariance reported in
item~1. In the generative direction the mechanism is far from degenerate: the attributes explain
only $R^2 = 0.14$ of the pixel variance, the remainder being digit identity and style, i.e.\ the
exogenous $\xi$.

\subsection{Constructing the observational dataset}

The FF targets the observational setting, in which each unit is seen under exactly one arm.
The training set is constructed accordingly: for each image we retain either the untreated
outcome ($T = 0$) or the treated one ($T = 1$), never both,
\begin{equation}
    Y_i^{\mathrm{obs}} = T_i \, Y_i(1) + (1 - T_i)\, Y_i(0).
\end{equation}
The model therefore never observes paired data. The held-out pairing $\{Y_i(0), Y_i(1)\}$ is
retained by the generator under evaluation-only keys and used exclusively for scoring.

\subsection{Engineering the confounding}

Treatment is assigned by a logistic propensity in the standardised thickness
$\tilde{A}_i = (A_i - \bar{A})/\mathrm{sd}(A)$:
\begin{equation}
    \pi_i \;=\; \Pr(T_i = 1 \mid A_i) \;=\; \sigma\!\left(\beta_0 + \beta_1 \tilde{A}_i\right),
    \qquad T_i \sim \mathrm{Bernoulli}(\pi_i).
    \label{eq:propensity}
\end{equation}
Here $\beta_1$ is the confounding strength: $\beta_1 = 0$ gives a randomised controlled trial,
and larger $\beta_1$ makes structurally thicker digits progressively more likely to be
treated. The intercept $\beta_0$ controls the marginal treated fraction and is $0$ throughout,
giving approximate $50/50$ balance. Because thickness also drives the untreated image,
$\beta_1 > 0$ opens a genuine back-door path and an unadjusted contrast is biased.

The value $\beta_1 = 1.2$ used below yields $\operatorname{corr}(A, T) \approx 0.45$ with
propensities spanning $[0.06, 0.999]$ --- strong confounding while retaining overlap. Both
extremes are reported per run, since a $\beta_1$ large enough to threaten positivity would
make the estimand non-identified for reasons unrelated to the FF.

\subsection{Constructing the effect}
\label{sec:effect-construction}

The modulation $c_{ik}$ is assembled from mean-zero \emph{rank scores}. For a unit-level
covariate $W$ with within-sample rank $R_i \in \{0, \ldots, n-1\}$, define the rank score
$u_i = (R_i + \tfrac12)/n$ and the modulator
\begin{equation}
    h_i \;\propto\; \varphi(u_i) - \frac{1}{n}\sum_{j=1}^{n} \varphi(u_j),
    \label{eq:modulator}
\end{equation}
rescaled so that $\max_i |h_i| = 1$. Three ingredients are built this way: $h^{A}$ from
thickness, $h^{B}$ from brightness, and $g_{ik}$ from the \emph{within-pixel} rank of the
unit's own untreated outcome $Y_i(0)_k$.

\paragraph{Exactness of the imposed ATE.}
Because Equation~\eqref{eq:modulator} subtracts the \emph{empirical} mean rather than an
analytic one, $\bar{h} = 0$ holds identically --- for any $n$, any data, and crucially for any
shaping function $\varphi$. Every modulation below is affine in such terms, so
\begin{equation}
    \mathrm{ATE}_k \;=\; \frac{1}{n}\sum_{i=1}^n \tau_{ik}
    \;=\; m_k \cdot \frac{1}{n}\sum_{i=1}^n c_{ik}
    \;=\; m_k .
    \label{eq:exact-ate}
\end{equation}
The imposed ATE is thus the effect map itself: a quantity the experimenter \emph{sets} rather
than one measured from the realised sample. The generator asserts
Equation~\eqref{eq:exact-ate} at build time; the realised discrepancy is $\sim 10^{-15}$,
i.e.\ floating-point rounding alone.

It is the centring, not any linearity, that delivers this. Consequently $\varphi$ may be
non-monotone or discontinuous, and the modulation may be given arbitrary spatial structure,
without disturbing the truth. We exploit both freedoms below. A second consequence of the
rank construction is that every modulator is bounded in $[-1,1]$, so the effect retains the
sign of the map whenever the coefficients sum below one; the realised fraction of sign flips
is reported per run.

\paragraph{Modes.} The modulation takes one of three forms, distinguished by what the effect is
permitted to depend on.
\begin{itemize}
    \item \textbf{Covariate-only.}
    $c_{ik} = 1 + a_A h^{A}_i + a_B h^{B}_i + a_{AB}\,\widetilde{h^{A}_i h^{B}_i}$,
    where the tilde denotes re-centring of the product. The effect depends on observed
    pretreatment covariates alone; the interaction makes the conditional effect surface
    non-additive. This is unambiguously treatment effect heterogeneity.

    \item \textbf{Outcome-coupled.} $c_{ik} = 1 + a_A h^{A}_i + b\, g_{ik}$. The second term is
    keyed on the unit's own $Y(0)$ rank, so it is a \emph{coupling between the potential
    outcomes} rather than covariate heterogeneity. Such a coupling is not identified from
    observational data: only its consequence for the marginal law of $Y(1)$ is. We therefore
    keep the two effects in separate experiments rather than conflating them.

    \item \textbf{Quantile-primitive.} The quantile-varying effect is specified directly as
    $\delta_k(u) = m_k\big(1 + b\,\psi(u)\big)$ and $Y(1)$ constructed from it, so that
    \begin{equation}
        Q^{(1)}_k(u) - Q^{(0)}_k(u) \;=\; \delta_k(u)
        \label{eq:quantile-primitive}
    \end{equation}
    holds exactly and the target is available in closed form. This requires the shift to be a
    function of $(k, u)$ alone, so no unit-level term is admitted in this mode; the generator
    additionally verifies that the treated margin remains monotone, without which
    Equation~\eqref{eq:quantile-primitive} would fail.
\end{itemize}

\paragraph{Shape of the conditional effect.} The function $\varphi$ in
Equation~\eqref{eq:modulator} determines how the conditional effect varies with the covariate.
Beyond the linear default we use a cubic (monotone, flat in the middle and steep in the
tails), a quadratic (U-shaped, so extreme units respond and typical ones do not), a sinusoid
(fully non-monotone), and a step (discontinuous, two groups). These provide conditional effect
functions that no linear-in-covariate estimator can match by construction.

\paragraph{Spatial structure.} Replacing the scalar coefficient $a_A$ by a per-pixel
coefficient $a_A + a_{S}\,\psi_k$, with $\psi$ a smooth spatial pattern (horizontal or vertical
gradient, diagonal, or radial) centred and scaled over the effect's support, makes units differ
in the \emph{shape} of their effect rather than only its magnitude: thick digits respond more
on one side of the image and thin digits on the other, while the average over units remains
exactly $m$. This is the spatially coherent heterogeneity a real lesion would exhibit, and by
Equation~\eqref{eq:exact-ate} it costs nothing in exactness.

\subsection{The experiment ladder}

\begin{table}[h]
\centering
\small
\begin{tabular}{llll}
\toprule
 & \textbf{Confounded} & \textbf{Effect depends on} & \textbf{ATE\,=\,ATT?} \\
\midrule
E1 & no  & --- (identical for all units)              & yes \\
E2 & yes & ---                                        & yes \\
E3 & yes & thickness \emph{and} own $Y(0)$ rank       & no  \\
E4 & yes & thickness, brightness, interaction         & no  \\
E5 & yes & outcome quantile $u$ only                  & no  \\
E6 & yes & as E4, plus a spatial gradient             & no  \\
\bottomrule
\end{tabular}
\caption{The six experiments as settings of the shared generator. E1 $\to$ E2 changes only the
assignment mechanism; E2 $\to$ E3--E6 changes only the effect; E4 $\to$ E6 adds only spatial
structure.}
\label{tab:ladder}
\end{table}

\paragraph{E1 --- randomised, homogeneous.} Treatment is a fair coin and every treated unit
receives the identical map. The FF has no confounding to undo and no heterogeneity to
represent, so E1 measures the method's floor: whatever error remains is attributable to the
flow's approximation of a $K$-dimensional joint outcome, not to the causal problem.

\paragraph{E2 --- confounded, homogeneous.} The effect is unchanged from E1; only the
assignment mechanism moves, so the E1 $\to$ E2 difference is attributable to confounding
alone. Because the effect is a pure location shift, $\mathrm{ATE} = \mathrm{ATT} =
\mathrm{ATC}$ exactly and the true quantile-resolved effect is exactly flat in $u$ for every
pixel. E2 is therefore also a specification test: any curvature a flexible margin reports is
spurious.

\paragraph{E3 --- confounded, outcome-coupled.} Adds both a covariate term and a term keyed on
the unit's own $Y(0)$ rank. The second is a coupling between the potential outcomes rather
than treatment effect heterogeneity, and we are careful about what this does and does not
affect. It does \emph{not} disturb identification: $T$ is generated from $A$ and an independent
uniform whatever form $\tau$ takes, so $T \perp (Y(0), Y(1)) \mid A$ continues to hold and the
ATE remains identified --- E3 is a valid ATE experiment. What it does mean is that $\tau$ is
not a function of pretreatment variables alone, so a conditional average treatment effect is
not well defined for that component, and the coupling itself is not recoverable from
observational data since only one potential outcome per unit is ever observed. E3 is therefore
retained as the mixed case and reported as such; E4 and E6 supply heterogeneity with no
outcome coupling, and E5 states the quantile dependence explicitly as a primitive rather than
obtaining it as a by-product of a coupling.

\paragraph{E4 --- confounded, covariate heterogeneity.} The effect varies with thickness,
brightness and their interaction, with no dependence on $Y(0)$ whatsoever. This is the clean
conditional-effect experiment.

\paragraph{E5 --- confounded, quantile effect.} The quantile-varying effect is the primitive,
so the estimation target is available in closed form rather than measured. This isolates
whether a flexible margin recovers an effect that varies across the outcome distribution.

\paragraph{E6 --- confounded, spatially coherent heterogeneity.} As E4, with the covariate's
influence additionally shaped across pixels, so different units receive differently-shaped
effects that average to the common map.

\paragraph{}
Two properties make E3--E6 strictly more informative than E1--E2. First, confounding
correlates $T$ with thickness while thickness modulates the effect, so $\mathrm{ATT} \ne
\mathrm{ATE} \ne \mathrm{ATC}$; all three are computed exactly from the retained potential
outcomes, and their comparison identifies \emph{which} estimand a method has in fact
recovered --- a distinction E1 and E2 are structurally blind to. Second, the modulation
structure is itself diagnostic: the matrix $[c_{ik}]$ has rank $1$ under E4, rank $2$ under
E6, and full rank under E3 and E5, so a covariate-driven effect is provably spanned by a
handful of unit-level scores whereas an outcome-coupled one is not.

\subsection{Estimands and evaluation targets}
\label{sec:targets}

The primary target is the per-pixel ATE of Equation~\eqref{eq:exact-ate}, scored against the
estimate $\hat{\tau} \in \R^K$ by mean absolute error, RMSE, maximum absolute error, and
spatial correlation. Errors are reported \emph{separately on and off the effect's support}:
the map is nonzero on only $12$ of $64$ pixels, so a pooled average is dominated by pixels
whose true effect is exactly zero, and "recovered the magnitude" and "kept the zeros at zero"
are distinct failure modes. $\mathrm{ATT}$ and $\mathrm{ATC}$ are computed on the corresponding
sub-samples and reported alongside.

For the quantile-resolved effect we distinguish two functionals, which coincide when the DGP
is rank-preserving and separate otherwise:
\begin{align}
    \tau^{\mathrm{paired}}_k(u) &= \E\big[\, \tau_{ik} \;\big|\; \mathrm{rank}(Y_i(0)_k) = u \,\big],
    \label{eq:tau-paired}\\
    \tau^{\mathrm{marg}}_k(u)   &= Q^{(1)}_k(u) - Q^{(0)}_k(u).
    \label{eq:tau-marginal}
\end{align}
Equation~\eqref{eq:tau-paired} is the average individual effect at outcome rank $u$;
Equation~\eqref{eq:tau-marginal} contrasts the two margins. A flow whose causal margin is
monotone in $u$ can represent the latter but not the former, so
\textbf{Equation~\eqref{eq:tau-marginal} is the estimation target}, with the gap between the
two reported as a property of the design. Both integrate over $u$ to the ATE. Under E5 the
target is additionally known in closed form via
Equation~\eqref{eq:quantile-primitive}.

\subsection{Estimator arms}
\label{sec:arms}

Each experiment is fitted under two causal-margin parameterisations, which return a
commensurable $\hat{\tau} \in \R^K$ and are therefore directly comparable.

\begin{itemize}
    \item \textbf{Location translation.} Treatment is masked from the margin flow, so the
    entire effect is carried by a per-pixel location shift. $\hat{\tau}$ is an exact model
    parameter, read directly off the fitted chain. Correctly specified wherever the effect is
    a pure shift; misspecified wherever it varies with $u$.

    \item \textbf{Flexible continuous.} Treatment enters the $K$-dimensional spline margin
    directly, so $\hat{\tau}$ is a Monte-Carlo functional of the fitted interventional
    distribution, estimated by paired common-random-number sampling of $Y \mid do(T=0)$ and
    $Y \mid do(T=1)$; this arm additionally returns $\hat{\tau}_k(u)$. Correctly specified
    throughout, so on E1--E2 it quantifies the cost of unneeded flexibility. It runs on either
    of two conditioner engines --- a MADE-masked MLP, or a causal-transformer conditioner
    following TarFlow --- which are compared at identical capacity.
\end{itemize}

Confounding is handled in both arms by conditioning on the stage-one covariate ranks $U_Z$,
not by weighting.

Optimisation is not neutral between these arms and we tune it per arm rather than imposing a
common setting. The location-translation arm carries an explicit additive parameter that must
travel from its initialisation to the true effect and under-trains at small learning rates,
while the transformer conditioner can diverge at large ones; a single shared rate silently
advantages whichever arm it happens to suit.

\subsection{The sampling noise floor}
\label{sec:floor}

Because the ATE is imposed exactly, the only irreducible error is the finite sample. We
calibrate it two ways. Under E1, where assignment is random, the unadjusted difference in
means is unbiased and its deviation from the truth is pure sampling noise: at $n = 5923$ and
$s = 8$ this is $0.064$ in maximum absolute error across pixels. Under E2--E6,
inverse-probability weighting by the \emph{true} propensity of
Equation~\eqref{eq:propensity} --- an oracle unavailable to any real estimator --- returns
$0.068$, statistically indistinguishable from the E1 floor.

This serves as both a sanity check and a target. The oracle result confirms each design is
identified, so any failure is attributable to the estimator rather than to a broken DGP; and
$\approx 0.065$ is the resolution limit at this sample size, so it, rather than zero, is the
value against which recovery should be judged. For reference, the unadjusted contrast is in
error by $0.75$ to $0.97$ across E2--E6 at the same scale, so the confounding to be undone is
an order of magnitude larger than the noise floor.

\subsection{Implementation}

The generator and the estimation script are self-contained and archive every run with its full
configuration, the git commit, library versions, the fitted scores, and the raw arrays needed
to reproduce every figure without refitting. A built-in correctness suite asserts the
invariants relied on above --- among them that the imposed ATE equals the effect map to
machine precision under every mode, shape and spatial basis; that the potential outcomes and
the observed mixture are consistent; that the homogeneous experiments have coincident
estimands and a flat quantile effect while the heterogeneous ones do not; that
Equation~\eqref{eq:quantile-primitive} holds under E5; and that each estimator arm returns a
finite per-pixel estimate of the correct dimension.

\paragraph{Extension: hidden confounding.} The experiments above confound through fully
observed covariates. Unobserved confounding can be introduced by correlating the treatment and
outcome ranks through a Gaussian copula with parameter $\rho$, following the standard FF
mechanism, under which the estimand remains exactly known while back-door adjustment ceases to
be valid. This is not exercised in the experiments reported here.
